\SourceLeaf{455}
\par
4) Man hat nun noch die Gleichung
\[
dZ=PdX+QdY
\]
\noindent zu integriren, wozu noch
\[
dZ^2+dX^2+dY^2=dz^2+dx^2+dy^2
\]
\noindent zu setzen ist.

\Needspace{9\baselineskip}
\begin{center}\textbf{BEMERKUNGEN.}\end{center}
\noindent\hspace*{1.2em}Wenn man die Coordinaten zweier correspondirender Punkte $P$ und $P'$ der auf einander abwickelbaren Flächen $S$ und $S'$ als Functionen zweier unabhängig veränderlicher Variabelen $p,q$ dargestellt denkt, so wird das Quadrat des Linienelements dieser Flächen durch den quadratischen Differentialausdruck
\[
dx^2+dy^2+dz^2=dx'^2+dy'^2+dz'^2=Edp^2+2Fdpdq+Gdq^2
\]
\noindent dargestellt. Betrachtet man gleichzeitig die zwei sphärischen Abbildungen dieser Flächen auf die Einheitskugel, bezeichnet durch $\xi,\eta,\zeta$ die Coordinaten des Bildes von $P$, durch $\xi',\eta',\zeta'$ die entsprechenden des Bildes von $P'$, und bestimmt die Quadrate der Linienelemente dieser Abbildungen aus den Gleichungen:
\[
\begin{array}{ll}
\xi=\sqrt{1-\zeta^2}\cos u,&\eta=\sqrt{1-\zeta^2}\sin u,\\[4pt]
\xi'=\sqrt{1-\zeta'^2}\cos u',&\eta'=\sqrt{1-\zeta'^2}\sin u',
\end{array}
\]
\noindent so erhält man:
\[
\begin{aligned}
d\xi^2+d\eta^2+d\zeta^2&=\frac{d\zeta^2}{1-\zeta^2}+(1-\zeta^2)du^2\\[5pt]
d\xi'^2+d\eta'^2+d\zeta'^2&=\frac{d\zeta'^2}{1-\zeta'^2}+(1-\zeta'^2)du'^2.
\end{aligned}
\]

\noindent Führt man die Grössen $\zeta,\zeta'$ als diejenigen Variabeln ein, durch welche die Lage der Punkte $P$ und $P'$ in $S$ und $S'$ bestimmt werden, so folgen die drei nachstehenden Gleichungen:
\[
dx^2+dy^2+dz^2=dx'^2+dy'^2+dz'^2=Ed\zeta^2+2Fd\zeta d\zeta'+Gd\zeta'^2
\]
\[
\begin{aligned}
d\xi^2+d\eta^2+d\zeta^2
&=\frac{d\zeta^2}{1-\zeta^2}+(1-\zeta^2)\left[\frac{\partial u}{\partial\zeta}d\zeta+\frac{\partial u}{\partial\zeta'}d\zeta'\right]^2\\
&=\mathfrak{E}d\zeta^2+2\mathfrak{F}d\zeta d\zeta'+\mathfrak{G}d\zeta'^2
\end{aligned}
\]
\[
\begin{aligned}
d\xi'^2+d\eta'^2+d\zeta'^2
&=\frac{d\zeta'^2}{1-\zeta'^2}+(1-\zeta'^2)\left[\frac{\partial u'}{\partial\zeta}d\zeta+\frac{\partial u'}{\partial\zeta'}d\zeta'\right]^2\\
&=\mathfrak{E}'d\zeta^2+2\mathfrak{F}'d\zeta d\zeta'+\mathfrak{G}'d\zeta'^2.
\end{aligned}
\]

\noindent Stellt $k$ den Werth des Krümmungsmasses im Punkte $P$ von $S$ oder im Punkte $P'$ von $S'$ dar, so gelten ferner die Beziehungen:
\[
\mathfrak{E}\mathfrak{G}-\mathfrak{F}^2=k^2(EG-F^2)=\mathfrak{E}'\mathfrak{G}'-\mathfrak{F}'^2,
\]

